\documentclass[10pt,a4paper]{amsart}
\usepackage[margin=19mm]{geometry}
\usepackage{amsmath,amssymb,mathtools}
\usepackage[scr=rsfs,cal=euler]{mathalfa}
\usepackage{xfrac}
\usepackage[unicode,hidelinks,bookmarks=false]{hyperref}
\newcommand{\ie}{i.e.\ }
\newcommand{\cf}{cf.\ }
\newcommand{\resp}{respectively\ }
\newcommand{\Subsection}{\subsection}
\providecommand{\nobreakdash}{\nobreak\hbox{-}}
\setlength{\emergencystretch}{2em}
\multlinegap0pt
\usepackage{bm}
\usepackage{bbm}
\usepackage{dsfont}
\usepackage{xspace}
\usepackage{cancel}
\usepackage{mathtools}
\usepackage{amsthm}
\makeatletter
\@ifundefined{theorem}{\newtheorem{theorem}{Theorem}}{}
\@ifundefined{definition}{\newtheorem{definition}[theorem]{Definition}}{}
\@ifundefined{lemma}{\newtheorem{lemma}[theorem]{Lemma}}{}
\makeatother
\newcommand{\R}{\mathbb R}
\title[A Riemannian viewpoint on the Amari-Cencov $\alpha$-connecti]{A Riemannian viewpoint on the Amari-Cencov $\alpha$-connections and Proudman-Johnson equations}
\author{Martin Bauer, Alice Le Brigant, Cy Maor}
\date{Source: arXiv:2508.00371v1 (CC BY 4.0)}
\begin{document}
\maketitle

\noindent\emph{Source and licence.} A Riemannian viewpoint on the Amari-Cencov $\alpha$-connections and Proudman-Johnson equations. Version arXiv:2508.00371v1 (CC BY 4.0), distributed under the Creative Commons Attribution 4.0 International licence, as recorded in the arXiv licence metadata for this exact version and shown on the abstract page https://arxiv.org/abs/2508.00371v1. Full rights notice: licenses/arxiv-2508.00371-CC-BY-4.0.txt.\par\smallskip

\noindent\textit{Editorial scope.} Counted unit: the Lemma whose source label is \texttt{lem:exponential} in the article (source0.tex lines 1095--1111, source label \texttt{lem:exponential}), delivered together with the article's own complete original proof of it (source0.tex lines 1113--1138, one \texttt{proof} environment of 26 lines). No mathematics is written, repaired or re-derived: statement and proof are the article's own text line for line. Required context retained uncounted: exp\_family (lines 1086--1092). Every definition, display and label of those lines is retained unchanged. Omitted from this package and not redistributed: 1--1085, 1093--1094, 1112--1112, 1139--1325. Nothing inside a retained excerpt is omitted or altered: every retained body is delivered exactly as the article's lines are written, so no omission receipt of any kind accompanies this package. Citations: the retained excerpts carry no citation command at all, so no bibliography entry of the article is transcribed and no reference is redistributed. Reference adaptation: none was needed and none was made; every reference inside the delivered unit resolves inside it, so no unresolved reference or citation is produced. Printed slips kept literal: no printed word, symbol, hypothesis or number of the counted statement or of its proof was altered. Layout and class: the article's own class is \texttt{amsart} with packages including mathtools, xcolor, amssymb, enumitem, geometry, inputenc, aeguill, hyperref, todonotes, xy, stmaryrd, appendix, multirow, ulem; the wrapper is the lane's pinned \texttt{amsart} editorial wrapper at 10pt on A4, which restates the theorem-like environments the retained text needs and adds the article's own macro definitions that the retained text needs (\texttt{\textbackslash R}), with extra wrapper packages (bm, bbm, dsfont, xspace, cancel, mathtools). The delivered numbering is the unit's own. Assets: no retained span contains a figure environment, a graphics-inclusion command, a PDF-page inclusion, a table or a TikZ picture; no image file is copied into this package, the package declares no asset dependency and no image file is redistributed with it. Licence: version arXiv:2508.00371v1 of this article is distributed under the Creative Commons Attribution 4.0 International licence, as recorded in the arXiv licence metadata for this exact version and shown on the abstract page https://arxiv.org/abs/2508.00371v1. A full rights notice is delivered at licenses/arxiv-2508.00371-CC-BY-4.0.txt. Characters: every retained excerpt is pure ASCII, so the delivered engine, XeLaTeX with its default Latin Modern text font, reports no missing character.
\begin{definition}[Exponential family]
A parametric model is an exponential family with natural parameter $\theta$ if the density is of the form
\begin{equation}\label{exp_family}
p(x,\theta)=\exp(h(x)+\theta\cdot\varphi(x)-F(\theta))
\end{equation}
for certain functions $h$, $\varphi$ on $\R^n$ and $F$ on $\Theta$.
\end{definition}

\begin{lemma}[Geometry of exponential families]\label{lem:exponential}
Let $\mathcal P_\Theta = \{p(x,\theta)\,dx: \theta\in\Theta\}$ be an exponential family, i.e., such that $p$ is of the form \eqref{exp_family}. Then for $1\leq i,j,k,\ell\leq d$:
\begin{itemize}
\item The Fisher--Rao metric is the hessian of the log-partition function $F$
\begin{equation}\label{exp_fisher_rao}
G_{ij}(\theta)=\partial_i\partial_j F(\theta).
\end{equation}
\item The Christoffel symbols of the $\alpha$-connections are given by
\begin{equation}\label{exp_alpha}
\Gamma^{(\alpha)}_{ij,k}=\frac{1-\alpha}{2}F_{ijk}.
\end{equation}
\item The curvature tensor is written
\begin{equation}\label{exp_alpha_curv}{R^{(\alpha)}}_{ijk}^\ell=\frac{1-\alpha^2}{4}G^{rm}G^{s\ell}\left(F_{ikm}F_{jrs}-F_{jkm}F_{irs}\right)
\end{equation}
\end{itemize}
Here we have used the notation $F_{ijk}:=\partial_i\partial_j\partial_kF$.
\end{lemma}

\begin{proof}
The expression for $G_{ij}$ follows from straightforward computation. Since the hessian of the log-density is independent of the sample variable and $\partial_k\ell(X,\theta)$ has zero expectation, we get
\begin{align*}
\Gamma^{(\alpha)}_{ij,k}&=-\partial_i\partial_jF(\theta)E_\theta\big[\partial_k\ell(X,\theta)\big]+\frac{1-\alpha}{2}E_\theta\big[\partial_i\ell(X,\theta)\partial_j\ell(X,\theta)\partial_k\ell(X,\theta)\big]\\
&=\frac{1-\alpha}{2}E_\theta\big[\partial_i\ell(X,\theta)\partial_j\ell(X,\theta)\partial_k\ell(X,\theta)\big].
\end{align*}
Denoting $\partial_if=f_i$, we notice that, since $\ell_{ij}=f_{ij}/f- f_{i}f_j/f^2$ and $\ell_{ij}$ is constant in the sample variable,
\begin{align*}
E_\theta\big[\ell_{ijk}\big]&=E_\theta\left[\partial_i\left(\frac{f_{jk}}{f}-\frac{f_jf_k}{f^2}\right)\right]=E_\theta\left[\frac{f_{ijk}}{f}-\frac{f_{jk}f_i}{f^2}-\frac{f_{ki}f_j}{f^2}-\frac{f_{ij}f_k}{f^2}+2\frac{f_if_jf_k}{f^3}\right]\\
&=0-E_\theta\big[\ell_{jk}\ell_i\big]-E_\theta\big[\ell_{ki}\ell_j\big]-E_\theta\big[\ell_{ij}\ell_k\big]-E_\theta\big[\ell_i\ell_j\ell_j\big]\\
&=-\ell_{jk}E_\theta\big[\ell_i\big]-\ell_{ki}E_\theta\big[\ell_j\big]-\ell_{ij}E_\theta\big[\ell_k\big]-E_\theta\big[\ell_i\ell_j\ell_j\big]=-E_\theta\big[\ell_i\ell_j\ell_j\big],
\end{align*}
and since $E_\theta\big[\ell_{ijk}\big]=-F_{ijk}$ we obtain $\Gamma_{ij,k}^{(\alpha)}=\frac{1-\alpha}{2}F_{ijm}$. Then the curvature tensor defined by $R^{(\alpha)}(\partial_i,\partial_j)\partial_k=:{R^{(\alpha)}}_{ijk}^\ell \partial_\ell$ is written
\begin{align*}
{R^{(\alpha)}}_{ijk}^\ell&=\partial_i\Gamma_{jk}^\ell -\partial_j\Gamma_{ik}^\ell+\Gamma_{im}^\ell\Gamma_{jk}^m-\Gamma_{jm}^\ell\Gamma_{ik}^m\\
&=\frac{1-\alpha}{2}\left(F_{ijkm}G^{m\ell}+F_{jkm}\partial_iG^{m\ell}-F_{ijkm}G^{m\ell}+F_{ikm}\partial_jG^{m\ell}\right)\\
&+\frac{(1-\alpha)^2}{4}\left(F_{jkr}g^{rm}F_{ims}G^{s\ell}-F_{ikr}G^{rm}F_{jms}G^{s\ell}\right).
\end{align*}
Then using the fact that $\partial_iG^{mr}G_{rs}=-G^{mr}\partial_iG_{rs}$ and so $\partial_iG^{m\ell}=-G^{mr}\partial_iG_{rs}G^{s\ell}$, we obtain
\begin{align*}
{R^{(\alpha)}}_{ijk}^\ell&=\frac{1-\alpha}{2}\left(F_{jkm}F_{irs}G^{rm}G^{s\ell}+F_{ikm}F_{jrs}G^{mr}G^{s\ell}\right)\\&\qquad+\frac{(1-\alpha)^2}{4}\left(F_{jkr}F_{ims}G^{rm}G^{s\ell}-F_{ikr}F_{jms}G^{rm}G^{s\ell}\right)\\
&=\frac{1-\alpha}{4}\bigg(-2F_{jkm}F_{irs}G^{rm}G^{s\ell}+F_{ikm}F_{jrs}G^{rm}G^{s\ell}\\&\qquad+(1-\alpha)F_{jkm}F_{irs}G^{rm}G^{s\ell}-(1-\alpha)F_{ikm}F_{jrs}G^{rm}G^{s\ell}\bigg)\\
&=\frac{1-\alpha^2}{4}G^{rm}G^{s\ell}\left(F_{ikm}F_{jrs}-F_{jkm}F_{irs}\right)
\end{align*}
as expected.
\end{proof}

\end{document}
